\subsection{An exhaustive split of the affine-linear coordinate}\label{sec:wsplit}

Write $h=a+bZ+cT+dZ^2+eZT+fT^2$ and
$\Delta=e^2-4df$. The following split is exhaustive and its branches
are disjoint. The last column distinguishes a proof given here from
an analytic obligation that has not yet been discharged in this draft.

\begin{center}\small
\begin{tabular}{@{}p{0.27\textwidth}p{0.38\textwidth}p{0.25\textwidth}@{}}
\toprule
Coefficient condition & Reduction & Status in this draft\\
\midrule
$\Delta\ne0$ & Laurent phase on $\C^*$ with logarithmic volume &
Uniform logarithmic cut estimate still required\\[3pt]
$\Delta=0$, $f\ne0$ & $h=a+bZ+cT+fT^2$ after a null rotation &
Proposition~\ref{prop:rankoneformal}\\[3pt]
$f=e=0$, $d\ne0$, $c\ne0$ &
$W(a+T+Z^2)+G$ after normalization &
Proposition~\ref{prop:mixedformal}\\[3pt]
$f=e=0$, $d\ne0$, $c=0$ & $Wh(Z)+G$ &
Proposition~\ref{prop:height}\\[3pt]
$d=e=f=0$, $c\ne0$ & $W(a+bZ+cT)+G$ &
Corollary~\ref{cor:linearwt}\\[3pt]
$d=e=f=c=0$ & $Wh(Z)+G$ &
Proposition~\ref{prop:height}\\
\bottomrule
\end{tabular}
\end{center}

Here is the algebra behind the reductions. The null rotation
\[
 Z=Z',\quad T=T'+\lambda Z',\quad
 W=W'-\lambda T'-\lambda^2Z'/2
\]
preserves $2ZW+T^2$. If $f\ne0$, choose $\lambda=-e/(2f)$.
When $\Delta=0$ the quadratic part of $h$ becomes just $fT'^2$.
If $f=e=0$ and $c\ne0$, choose $\lambda=-b/c$ to remove the
linear $Z$ term without changing $dZ^2$. If also $d\ne0$, the
Gaussian scaling $Z=\kappa Z'$, $W=W'/\kappa$, followed by a
nonzero scaling of $P$, makes the $T$ and $Z^2$ coefficients both
one; take $\kappa^2=c/d$. All additional terms enter a polynomial
$G$ of degree at most three. If $\Delta=0$ and $f=0$, then $e=0$,
which proves that the table omits no case.

\subsection{What the remaining analytic estimates must actually say}
\label{sec:analyticgates}

An integral representation is not an asymptotic estimate. The needed
estimates can be made elementary once an appropriate flat cut basis
has been constructed. The uniformity and completeness of that basis
are part of the proof, not consequences of a few saddle calculations.

\begin{lemma}[A conditional shrinking estimate]\label{lem:conditionalshrink}
Fix a ray $t=Re^{i\theta}$. Suppose the relevant zero-exponential
periods of $a(t)e^{\varphi(t,x)}dx$ have a spanning flat cut family
with fixed coefficients and the following properties:
\begin{enumerate}
\item On the entire phase band $|\varphi|\leq K\log R$, the lifted
cuts lie in a bounded number of disks of radius $\delta(R)$, where
$\delta(R)\to0$.
\item They are pieces of real algebraic curves of bounded degree,
with a bounded number of pieces, and $\operatorname{Re}\varphi\leq C$
there, uniformly in $R$.
\item The tails beyond the band $|\varphi|\leq K\log R$ are $o(1)$,
uniformly for the chosen flat generators; $a(t)$ is bounded.
\end{enumerate}
Then every such scalar period tends to zero.
\end{lemma}
\begin{proof}
A real algebraic curve of degree at most $D$ has length at most
$\pi D\rho$ in a disk of radius $\rho$. This follows by averaging
the number of intersections with lines and applying the one-dimensional
coarea formula; line components satisfy the same bound.
Thus the bounded pieces have total length $O(\delta(R))$ and bounded
integrand. Their integrals tend to zero. Add the uniformly negligible
tails.
\end{proof}

For a stable polynomial degree, $\varphi=t^p(V+O(t^{-1}))$, the
equation $\varphi=w$ on a logarithmic band forces its bounded roots
into disks of radius $O((\log R/R)^{1/\deg V})$ around the zeros of
$V$. This explains the shrinking scale. It does not by itself prove
that those pieces span the entire zero-exponential part, nor supply
the uniform tail estimate. In particular, for $p>1$ the lower phase
terms can still diverge, so they cannot simply be treated as a bounded
analytic amplitude.

\begin{lemma}[A conditional logarithmic estimate]\label{lem:conditionallog}
Suppose a spanning flat cut family for the zero-exponential part of
$e^{\varphi(t,x)}dx/x$ satisfies the bounded-degree and bounded-real-part
conditions above in an annulus $R^{-M}\leq|x|\leq R^M$. Assume its
end detours have bounded logarithmic length and its descending tails
have uniformly bounded integrals. Then its scalar periods are
$O(1+\log R)$. After multiplication by $1/t$, they tend to zero.
\end{lemma}
\begin{proof}
In $r\leq|x|\leq2r$, the curve-length estimate gives
$\int|dx|/|x|\leq2\pi D$. There are $O(\log R)$ dyadic annuli.
The integrand is bounded by $e^C$; the end detours and tails add
bounded quantities. This proves the estimate.
\end{proof}

The unresolved step in applying these two lemmas is the construction
of one controlled \emph{flat} family with the claimed spanning
property, including regular puncture classes, coincident critical
values and degenerate critical points. Hien--Roucairol gives the
perfect pairing and representation; it does not assert these uniform
cut properties. They must be proved for the particular rational
phases below, or replaced by a complete formal-frame construction.

\subsection{The logarithmic branch, including its removed component}

For $\Delta\ne0$, $f\ne0$, choose $\eta^2=\Delta$ and set
\[
 \rho=\frac{e+\eta}{2f},\quad \sigma=\frac{e-\eta}{2f},\quad
 u=\frac{\rho c-b+1/t}{\eta},\quad
 v=\frac{b-1/t-\sigma c}{\eta},\quad k=uv-a/f.
\]
On the circular constraint $z=th(z,T)$, put
$x=T+\rho z+u$, so that $T+\sigma z+v=k/x$. Then
\[
 z=\frac f\eta(x-k/x-u+v),\quad T=x-u-\rho z,
 \qquad
 \frac{dT}{1-th_Z}=\frac{dx}{t\eta x}.
\]
The coefficient of $t^{-2}$ in $k$ is $-1/\eta^2$, so the conic
is generically smooth. The exponent $tG(z,T)-T^2/2$ is Laurent
polynomial in $x$. On the formal Gaussian branch,
$x=1/(\eta t)+O(1)$, which makes every Laurent power legitimate.

If $f=0$, then $e\ne0$. Put $z_0=-c/e$ and $x=z-z_0$. With
$A(z)=a+bz+dz^2$, the constraint gives
\[
 T=\frac{z/t-A(z)}{ex},\qquad
 \frac{dT}{1-th_Z}=\frac{dx}{tex}.
\]
This localization is valid if $c\ne0$, or if $c=0$ and $a\ne0$.
If $a=c=0$, the chosen fixed point is instead $g=0$ identically;
inverting $g$ would discard the very branch being tested. In that
case the exact unit series is
\[
 \E_T\left[\frac{e^{tG(0,T)}}{1-t(b+eT)}\right].
\]
Use $x=T-(1/t-b)/e$, now a valid Laurent unit. Its volume is
$-dx/(tex)$ and its exponent is polynomial in $x$. These formulas
cover both the smooth conic and the removed-component branch.
They reduce the exclusion to the precise logarithmic estimate just
stated; the formulas alone are not a proof of that estimate's
spanning and uniformity hypotheses.

\subsection{Why the polynomial branches no longer need a cut estimate}

Propositions~\ref{prop:rankoneformal} and~\ref{prop:mixedformal}
apply the same leading-coefficient criterion to their exact formal
Gaussian functionals. Their complete stationary-ball frames retain
all critical multiplicities and escaping roots. Consequently neither
proposition depends on the conditional shrinking lemma above.
The Laurent phase with logarithmic volume is different: a ball
containing the puncture may contribute logarithmic extensions.
The polynomial theorem does not cover that situation, and the
logarithmic obligation remains open in this draft.

\subsection{Adversarial tests for the analytic argument}

The following exact examples prevent several tempting shortcuts.
For positive real $t$,
\[
 \frac1t\int_{\mathbb R}e^{-x^2/(2t^k)}\,dx
   =\sqrt{2\pi}\,t^{k/2-1}.
\]
The unique critical point is bounded and has zero action for every
$k$, yet the normalized period decays for $k=1$, is constant for
$k=2$, and grows for $k=3$. A critical-point count or a zero-action
label alone gives no decay estimate.

For the fixed phase $tx^k$, $k>2$, the critical point is degenerate
and the decay scale is $t^{-1/k}$, not a Gaussian Hessian scale.
The fixed quotient in Theorem~\ref{thm:stationaryballs}
retains its full multiplicity $k-1$.
For $tV_nx^n-x^4/2$, $1\leq n\leq3$, the $n-1$ bounded and
$4-n$ escaping critical dimensions total three. Counting only
bounded points loses actual rank; the stationary-ball construction
retains all roots and tests its full matrix using polynomial separation.
Finally, $t^{-3/2}=r^3$ under $r=t^{-1/2}$: a claim about a
noninteger exponent before ramification must not be reused after
ramification. Lemma~\ref{lem:positiveformal} uses positive seed order
and an actual full frame instead.

\subsection{The geometric case split still required for frame extraction}

For a cubic of degree exactly three, either the frame test
\eqref{eq:frame} succeeds, or it does not. In the latter case,
the projective curve $F_3=0$ belongs to exactly one of these types:
\begin{center}\small
\begin{tabular}{@{}p{0.33\textwidth}p{0.56\textwidth}@{}}
\toprule
Leading curve & Required no-frame exclusion\\
\midrule
Smooth cubic & All finite critical blocks; no contribution from infinity\\
Irreducible nodal cubic & Isolated finite blocks and critical-line degenerations\\
Irreducible cuspidal cubic & Smooth parabola, separated lines and double-line critical profiles\\
Line and smooth conic, transverse & Both ordinary nodes and all affine escapes\\
Line and smooth conic, tangent & Tangency, degenerate finite blocks and escaping profiles\\
Three nonconcurrent lines & Finite critical sets, critical conics and axes\\
Three concurrent distinct lines & Constant and nonconstant transverse linear coefficient\\
Double line and a distinct line & Isotropic and nonisotropic repeated-line directions\\
Triple line & Isotropic and nonisotropic repeated-line directions\\
\bottomrule
\end{tabular}
\end{center}
The list is exhaustive by factorization of a plane cubic over $\C$.
It is a list of proof obligations, not a declaration that the exclusions
have been reproved here. For isolated leading singular points, a null
point with $F_2=0$ goes directly to the affine-linear-coordinate table;
otherwise the lower Gaussian quadratic and its cross terms must be
retained in the indicated exclusion. In the repeated-line cases the
singular locus is a whole projective line, and the frame test must be
applied to all its null directions, not to an arbitrarily selected point.

Thus the routing is now explicit. Completing frame extraction still
requires a proof for every no-frame row, including a full rank or
homology-completeness argument. The conditional estimates above say
exactly what an analytic proof must establish. They do not promote
finite algebra checks into that missing argument.
